\clearpage

\section {Introduction }
\label{sec:intro}

In this technical note, 
we report the results of a 3-flavour $\nu_{\mu}$-disappearance analysis,
performed by the Valencia-Lancaster-Oxford-Rutherford ({\em VaLOR}) group,        
on the combined Run 1+2+3+4a (6.3933$\times$10$^{20}$ POT) dataset.
This dataset was collected in four run periods:
in January - June 2010 (Run 1, 0.323$\times$10$^{20}$ POT),
in November 2010 - March 2011 (Run 2, 1.108$\times$10$^{20}$ POT) and,
following the recovery work from the damage of the 2011 Tohoku earthquake, 
in March - June 2012 (Run 3, 1.579$\times$10$^{20}$ POT) and
in October 2012 - 11th April 2013 (Run 4a, 3.3833$\times$10$^{20}$ POT).
This analysis is an update of the work presented in Refs.\ \cite{OurRun1DisappAnalysis}, \cite{OurRun12DisappAnalysis}, 
\cite{OurRun12DisappAnalysis3Flavour} and \cite{OurRun123DisappAnalysis3Flavour} with improved methods and an enlarged dataset. 

This note is meant to be read in conjunction with Ref.\ \cite{OurRun123DisappAnalysis3Flavour} since this analysis uses the same 
methods as our Run 1+2+3 analysis described in that reference,
with only a few of the analysis inputs being updated.
A summary of these updates is given in Sec.\ \ref{sec:changes}.

In this analysis, the oscillation parameters are determined
by fitting the predicted reconstructed energy spectrum
of 1-ring $\mu$-like events at SuperK to the observed spectrum.
The method of calculation of this predicted spectrum is presented in Sec.\ 2 of Ref.\ \cite{OurRun123DisappAnalysis3Flavour}.
This analysis uses flux tuning based on recent NA61 data \cite{FluxTuning}\cite{NA61T2K}. 
It also uses an extrapolation from ND280 to SuperK, which is implemented via a fit of flux and cross-section parameters to
ND280 measurements of muon momentum-angle distributions of $\nu_\mu$
CC-0-pion-like, CC-1-pion-like and CC-Other-like events whose 
vertices are in FGD1 \cite{BANFF2013}.
Updated nominal predictions of the reconstructed energy spectrum of single-ring $\mu$-like events in SuperK are presented for Run 1+2+3+4a 
in Sec.\ \ref{sec:spectrum_prediction}; the changes to the analysis summarized in Sec.\ \ref{sec:changes} are incorporated 
into these predictions. 

We consider a total of 44 energy scale, neutrino flux, neutrino cross section,
efficiency, initial-state nuclear environment and final-state rescattering,
and oscillation parameter systematics. 
These systematic parameters and their effects 
on the extrapolated single-ring $\mu$-like event reconstructed energy spectrum are described in Sec.\ 3
of Ref.\ \cite{OurRun123DisappAnalysis3Flavour}.

The oscillation analysis is performed using neutrino oscillation probabilities computed in a 3-flavour
framework including constant-density matter effects. 
Details of our oscillation probability calculations, estimates of the numerical accuracy of our
calculations and comparisons with calculations used by the alternative T2K oscillation analyses 
are given in Appendices A and B of Ref. \cite{OurRun123DisappAnalysis3Flavour}.

As in our previous analyses, the oscillation fit is made using the likelihood-ratio method \cite{StatAnalysisCowan}.
The mathematical details of this method are given in Sec.\ 3.2 in Ref.\ \cite{OurRun1DisappAnalysis}.
In this Run 1+2+3+4a analysis, we consider only a fit which allows all 44 systematic parameters to float,
as this was the version of our fit which was preferred by the collaboration during the Run 1+2+3 analysis.
The oscillation fitting method is presented in detail in Sec.\ 4.1 in Ref.\ \cite{OurRun123DisappAnalysis3Flavour}, 
while the chosen goodness-of-fit test, suitable for the current statistics, is described in Sec.\ 4.2 of  Ref.\ \cite{OurRun123DisappAnalysis3Flavour}.
The performance of the fitter was tested in detail for the Run 1+2+3+4a fit using ensembles of toy MC experiments, and
a summary of the results of these tests is presented in Sec.\ \ref{sec:sensitivity}.

The confidence regions were constructed using the Feldman-Cousins method \cite{FeldmanCousins} as in our previous analyses. 
Systematic errors were included in the oscillation contours using the 
Cousins-Highland method (Bayesian integration) \cite{CousinsHighland}.
Contours constructed with the constant-$\Delta\chi^{2}$ method are also shown for both fits for comparison.
Both of these methods are summarized in Sec.\ 4.5 of Ref.\ \cite{OurRun123DisappAnalysis3Flavour}, and
the exact methodology used for incorporating systematic errors is explained in detail.

% The expected T2K sensitivity for Run 1+2+3+4a was studied by making average confidence contours using the complete VaLOR analysis 
% for an integrated exposure of 6.3933$\times$10$^{20}$ POT. This was done for a grid of oscillation points 
% around the 2012 MINOS \cite{MINOS2012} 90\% CL region, and 
% the results are presented in Sec.\ \ref{sec:sensitivity}.

The results of the measurement of sin$^{2}\theta_{23}$ and $|\Delta m^{2}_{32}|$ 
using the Run 1+2+3+4a dataset are given in Sec.\ \ref{sec:results_bestfit} of this note, 
while the goodness-of-fit tests are presented in Sec.\ \ref{sec:results_gof}.
The best-fit spectra are shown for each of normal and inverted mass hierarchy assumptions for the Run 1+2+3+4a dataset
%Run 1+2+3, Run 4a and the combined Run 1+2+3+4a datasets
in Sec.\ \ref{sec:results_bestfit_spectra}.
The Feldman-Cousins and constant-$\Delta\chi^{2}$ confidence regions are shown in Sec.\ \ref{sec:results_contours}. 
Comparisons of our results obtained between different methods for constructing confidence regions are presented in Sec.\ \ref{sec:results_contour_comparisons}.
In that section, we also show comparisons of our T2K Run 1+2+3+4a results with results from MINOS and SuperK.
%as well as with the expected T2K sensitivity for an integrated exposure of 6.3933$\times$10$^{20}$ POT.

Finally, a summary is given in Sec.\ \ref{sec:summary}.  

